\documentclass[12pt,a4paper]{article}

\usepackage[a4paper,margin=18mm]{geometry}
\usepackage[T2A]{fontenc}
\usepackage[utf8]{inputenc}
\usepackage[russian]{babel}
\usepackage{amsmath,amssymb,mathtools}
\usepackage{array,booktabs,longtable}
\usepackage{xcolor}
\usepackage{hyperref}
\usepackage{tikz}

\hypersetup{
  colorlinks=true,
  linkcolor=blue!55!black,
  urlcolor=blue!55!black
}

\setlength{\parindent}{0pt}
\setlength{\parskip}{0.55em}
\setlength{\emergencystretch}{2em}
\renewcommand{\arraystretch}{1.25}
\newcolumntype{L}[1]{>{\raggedright\arraybackslash}p{#1}}

\begin{document}

% -------------------- Титульный лист --------------------
\begin{titlepage}
\centering
\vspace*{0.22\textheight}

{\LARGE\bfseries Ревью домашней работы\par}
\vspace{2.2cm}

{\large Автор: Лёвин Артём Александрович\par}
\vspace{0.8cm}

{\large 9 сентября 2026 года\par}

\vfill
\begin{tikzpicture}
  \draw[blue!35!black, line width=0.8pt]
    (0,0) .. controls (2.0,0.45) and (4.0,-0.45) .. (6.0,0)
          .. controls (8.0,0.45) and (10.0,-0.45) .. (12.0,0);
\end{tikzpicture}
\vspace{0.8cm}
\end{titlepage}

% -------------------- Сводная таблица --------------------
\newpage
\section*{Сводная таблица ошибок}

{\small
\setlength{\tabcolsep}{3pt}
\begin{tabular}{@{}L{2.2cm}L{3.0cm}L{3.9cm}L{2.5cm}L{2.9cm}@{}}
\toprule
\textbf{№ задания} & \textbf{Тема} & \textbf{Суть ошибки} & \textbf{Ваш ответ} & \textbf{Правильный ответ} \\
\midrule
\hyperref[task:3-2]{Упр. 3, №2} &
Вычитание десятичных дробей &
Получен неверный результат при поразрядном вычитании десятичных дробей; требуется точно выровнять разряды и корректно выполнить заимствование. &
$7{,}285$ &
$7{,}305$ \\
\bottomrule
\end{tabular}
}

% -------------------- Разделитель --------------------
\newpage
\thispagestyle{empty}
\vspace*{\fill}
\begin{center}
{\LARGE\bfseries Разбор ошибок}
\end{center}
\vspace*{\fill}

% -------------------- Задание 3.2 --------------------
\newpage
\phantomsection
\label{task:3-2}
\section*{Упражнение 3, №2}

\textbf{Текст задания.}
Вычислите:
\[
10{,}04-2{,}735.
\]

\textbf{Ваше решение.}
Достоверно читается итоговая запись
\[
10{,}04-2{,}735=7{,}285.
\]
Промежуточная запись столбиком частично зачёркнута и переработана, поэтому её отдельные исправления не восстанавливаются по догадке. Для проверки достаточно читаемых исходного выражения и итогового результата.

\textbf{Ваш ответ.}
\[
7{,}285.
\]

\textcolor{red}{\textbf{Ошибка:}} неверно выполнено вычитание десятичных дробей по разрядам.

\textbf{Где именно возникает ошибка.}
Последний достоверно правильный фрагмент Вашей записи --- само исходное действие
\[
10{,}04-2{,}735.
\]
Первый достоверно неправильный шаг --- переход к равенству
\[
10{,}04-2{,}735=7{,}285.
\]
Это равенство неверно. Конкретный зачёркнутый промежуточный перенос в записи столбиком нельзя надёжно прочитать, поэтому причина не приписывается одному неразборчивому штриху. Ошибка фиксируется на уровне результата поразрядного вычитания.

\textbf{Почему этот шаг неверен.}
При сложении и вычитании десятичных дробей числа необходимо записывать так, чтобы запятые стояли друг под другом. Тогда единицы вычитаются из единиц, десятые --- из десятых, сотые --- из сотых, тысячные --- из тысячных.

У числа $10{,}04$ можно без изменения его значения дописать справа ноль:
\[
10{,}04=10{,}040.
\]
После этого оба числа имеют по три знака после запятой:
\[
\begin{array}{r}
10{,}040\\[-2mm]
-\;2{,}735\\
\hline
\end{array}
\]
Если выполнять вычитание строго по разрядам, получается $7{,}305$, а не $7{,}285$.

\textcolor{blue!60!black}{\textbf{Ключевая идея:}}
перед вычитанием десятичных дробей нужно выровнять количество знаков после запятой, а затем выполнять обычное поразрядное вычитание справа налево, аккуратно учитывая каждое заимствование.

\textbf{Исправление от последнего правильного шага.}
Начинаем с исходного действия:
\[
10{,}04-2{,}735.
\]
Дописываем ноль справа у первой дроби:
\[
10{,}040-2{,}735.
\]
Записываем числа столбиком, совмещая запятые:
\[
\begin{array}{r}
10{,}040\\[-2mm]
-\;2{,}735\\
\hline
7{,}305
\end{array}
\]
Следовательно,
\[
10{,}04-2{,}735=7{,}305.
\]

\textbf{Полное правильное решение.}
Рассмотрим вычитание по разрядам подробно.

\begin{enumerate}
  \item Сначала представим первую дробь с тремя знаками после запятой:
  \[
  10{,}04=10{,}040.
  \]
  Дописывание нуля справа в дробной части не изменяет значение числа.

  \item Вычитаем тысячные. В числе $10{,}040$ в разряде тысячных стоит $0$, а нужно вычесть $5$. Заимствуем одну сотую из разряда сотых: $4$ сотых превращаются в $3$ сотых, а в разряде тысячных получаем $10$ тысячных. Поэтому
  \[
  10-5=5.
  \]

  \item Вычитаем сотые. После предыдущего заимствования осталось $3$ сотых. Поэтому
  \[
  3-3=0.
  \]

  \item Вычитаем десятые. В уменьшаемом стоит $0$ десятых, а вычесть нужно $7$ десятых. Непосредственно взять одну единицу из разряда единиц нельзя, потому что в числе $10$ в разряде единиц стоит $0$. Поэтому сначала одну десятку превращаем в $10$ единиц. Затем одну из этих единиц превращаем в $10$ десятых. После этого в разряде единиц остаётся $9$, а в разряде десятых становится $10$. Получаем
  \[
  10-7=3.
  \]

  \item Вычитаем единицы:
  \[
  9-2=7.
  \]

  \item В разряде десятков после заимствования остаётся $0$. Собираем полученные цифры по разрядам:
  \[
  7{,}305.
  \]
\end{enumerate}

Итак,
\[
\boxed{10{,}04-2{,}735=7{,}305}.
\]

\textbf{Проверка.}
Проверим вычитание обратным действием --- сложением разности и вычитаемого:
\[
7{,}305+2{,}735=10{,}040=10{,}04.
\]
Получили исходное уменьшаемое, значит результат $7{,}305$ верен.

Дополнительная оценка здравого смысла тоже согласуется с ответом: из числа, немного большего $10$, вычитается число, немного меньшее $3$, поэтому результат должен быть немного больше $7$. Число $7{,}305$ этому соответствует.

\textcolor{teal!60!black}{\textbf{Итог:}}
в этом примере правильный ответ --- $7{,}305$. Чтобы избежать такой ошибки, при вычитании десятичных дробей сначала выравнивайте запятые и при необходимости дописывайте нули справа, после чего выполняйте заимствования строго по разрядам.

\textbf{Задача на закрепление.}
Вычислите:
\[
12{,}3-4{,}875.
\]

\end{document}
